% META: {"id": "TNSI-LANG-EVAL-B", "chapitre": "TNSI-LANGAGES-ET-PROGRAMMATION", "type_objet": "evaluation", "version": "B", "capacites": ["T-LANG-01B", "T-LANG-03A", "T-LANG-03B", "T-LANG-03C", "T-LANG-04A", "T-LANG-04B", "T-LANG-05"], "mode_creation": "ex_nihilo", "sources_inspiration": [], "duree_min": 45, "status": "generated"}
\section*{Évaluation B --- Modules, paradigmes et débogage}
\textit{Durée : 45 minutes. Ordinateur avec interpréteur Python autorisé.}

\baremeIndicatif{Ex. 1 : 6 pts (module documenté) ; Ex. 2 : 8 pts (paradigmes, débogage) ; Ex. 3 : 4 pts (API et documentation) ; Ex. 4 : 2 pts (calculabilité, choix de paradigme)}

\bigskip

\textbf{Exercice 1} --- \textit{Créer un module documenté} \hfill (6 points)

\begin{enumerate}
  \item (4 pts) Écrire une fonction \lstinline{est_premier(n)}, destinée à un module \lstinline{arithmetique.py}, qui renvoie \lstinline{True} si l'entier \lstinline{n} (supérieur à $1$) est premier, \lstinline{False} sinon. Ajouter une docstring décrivant son rôle, son paramètre et sa valeur de retour.
  \item (2 pts) Tester la fonction sur $n=7$ (premier) et $n=8$ (non premier).
\end{enumerate}

% BEGIN-VERIFY
% def est_premier(n):
%     """Indique si l'entier n (n > 1) est premier."""
%     for d in range(2, n):
%         if n % d == 0:
%             return False
%     return True
% assert est_premier(7) is True
% assert est_premier(8) is False
% assert est_premier(2) is True
% assert est_premier.__doc__ is not None
% END-VERIFY

\bigskip

\textbf{Exercice 2} --- \textit{Paradigmes et débogage} \hfill (8 points)

\begin{enumerate}
  \item (3 pts) Réécrire en style fonctionnel (une liste en compréhension) la fonction impérative suivante :
\begin{python}
def carres_positifs_imperatif(valeurs):
    resultat = []
    for v in valeurs:
        if v > 0:
            resultat.append(v ** 2)
    return resultat
\end{python}
  \item (3 pts) Le code suivant contient un bug de nommage :
\begin{python}
def calculer_moyenne(notes):
    sum = 0
    for n in notes:
        sum += n
    return sum / len(notes)
\end{python}
    Quelle fonction native est masquée ? Ce code fonctionne-t-il malgré tout pour calculer une moyenne simple ? Justifier.
  \item (2 pts) Le code suivant compare deux résultats de calculs flottants avec \lstinline{==}. Que renvoie-t-il, pourquoi est-ce risqué de comparer ainsi, et comment corriger ?
\begin{python}
resultat = 0.3 - 0.2
print(resultat == 0.1)
\end{python}
\end{enumerate}

% BEGIN-VERIFY
% def carres_positifs_fonctionnel(valeurs):
%     return [v ** 2 for v in valeurs if v > 0]
% def carres_positifs_imperatif(valeurs):
%     resultat = []
%     for v in valeurs:
%         if v > 0:
%             resultat.append(v ** 2)
%     return resultat
% donnees = [-2, 3, -1, 4]
% assert carres_positifs_fonctionnel(donnees) == carres_positifs_imperatif(donnees)
% assert carres_positifs_fonctionnel(donnees) == [9, 16]
%
% def calculer_moyenne(notes):
%     sum = 0
%     for n in notes:
%         sum += n
%     return sum / len(notes)
% assert calculer_moyenne([10, 20, 30]) == 20.0
%
% resultat = 0.3 - 0.2
% assert resultat != 0.1
% assert abs(resultat - 0.1) < 1e-9
% END-VERIFY

\bigskip

\textbf{Exercice 3} --- \textit{Utiliser une API et sa documentation} \hfill (4 points)

On utilise le module standard \lstinline{statistics}, dont la documentation indique :

\begin{console}
statistics.mean(data)      -> moyenne arithmetique ; leve StatisticsError si data est vide
statistics.median(data)    -> mediane ; pour un nombre pair de valeurs, moyenne des deux centrales
statistics.pstdev(data)    -> ecart-type de population (diviseur n)
statistics.stdev(data)     -> ecart-type d'echantillon (diviseur n-1) ; exige au moins 2 valeurs
\end{console}

\begin{enumerate}
  \item (2 pts) Ecrire le code qui, pour la serie \lstinline{[4, 8, 6, 10]}, affiche la
  moyenne et la mediane en utilisant cette bibliotheque plutot qu'en les recalculant.
  \item (1 pt) La documentation distingue \lstinline{pstdev} et \lstinline{stdev}. Laquelle
  faut-il choisir pour decrire exactement les quatre valeurs mesurees, et pourquoi ?
  \item (1 pt) Que se passe-t-il si l'on appelle \lstinline{stdev([5])} ? Citer la phrase de
  la documentation qui permet de le prevoir sans executer le code.
\end{enumerate}

% BEGIN-VERIFY
% import statistics
% serie = [4, 8, 6, 10]
% assert statistics.mean(serie) == 7
% # Nombre pair de valeurs : la mediane est la moyenne des deux centrales, 6 et 8.
% assert statistics.median(serie) == 7
% assert sorted(serie)[1:3] == [6, 8]
% # pstdev decrit la population observee, stdev estime une population plus large.
% assert statistics.pstdev(serie) != statistics.stdev(serie)
% assert statistics.pstdev(serie) < statistics.stdev(serie)
% assert abs(statistics.pstdev(serie) - 5 ** 0.5) < 1e-12
% # La documentation annonce l'erreur : stdev exige au moins deux valeurs.
% try:
%     statistics.stdev([5])
% except statistics.StatisticsError:
%     pass
% else:
%     raise AssertionError("stdev sur une seule valeur aurait du echouer")
% # mean sur une serie vide leve la meme famille d'erreur.
% try:
%     statistics.mean([])
% except statistics.StatisticsError:
%     pass
% else:
%     raise AssertionError("mean sur une serie vide aurait du echouer")
% END-VERIFY

\bigskip

\textbf{Exercice 4} --- \textit{Calculabilite et choix d'un paradigme} \hfill (2 points)

\begin{enumerate}
  \item (1 pt) Un eleve affirme : \og{} Ce probleme n'est pas calculable en Python, il
  faudrait le programmer en C. \fg{} Expliquer pourquoi cette phrase confond deux choses
  differentes.
  \item (1 pt) Pour ecrire l'interface graphique d'une application, puis pour transformer une
  longue liste de mesures en une liste filtree et transformee, quel paradigme choisiriez-vous
  dans chaque cas ? Justifier en une phrase par cas.
\end{enumerate}

% BEGIN-VERIFY
% # La calculabilite ne depend pas du langage : les deux ecritures, imperative
% # et fonctionnelle, calculent la meme chose, et un langage en vaut un autre
% # des lors qu'il est Turing-complet.
% def puissance_imperative(base, exposant):
%     resultat = 1
%     for _ in range(exposant):
%         resultat *= base
%     return resultat
%
% def puissance_recursive(base, exposant):
%     if exposant == 0:
%         return 1
%     return base * puissance_recursive(base, exposant - 1)
%
% for base in range(1, 5):
%     for exposant in range(0, 6):
%         attendu = base ** exposant
%         assert puissance_imperative(base, exposant) == attendu
%         assert puissance_recursive(base, exposant) == attendu
% # Le paradigme change l'ecriture, pas le resultat.
% mesures = [12, -3, 7, 0, 25]
% imperatif = []
% for m in mesures:
%     if m > 0:
%         imperatif.append(m * 2)
% fonctionnel = [m * 2 for m in mesures if m > 0]
% assert imperatif == fonctionnel == [24, 14, 50]
% END-VERIFY
